% Problem record op_903c729ed6eb9ce6. This TeX file is derived from
% database/problems_json/op_903c729ed6eb9ce6.json by scripts/sync-tex.mjs; edit the JSON record
% and rerun the script rather than editing this file.
\section{Average-case additive hardness of squared complex Gaussian hafnians}
\label{sec:op_903c729ed6eb9ce6}

\paragraph{Problem.}
Is additive estimation of squared complex Gaussian hafnians $\#\mathrm P$-hard under randomized polynomial-time oracle reductions?
For $n\geq1$, let $X=X^T\in\mathbb C^{2n\times2n}$ have zero diagonal.
Its entries above the diagonal are independent with density $\pi^{-1}e^{-|z|^2}$.
Let $\mathcal M_{2n}$ be the perfect matchings of $\{1,\ldots,2n\}$ and define
\begin{equation}
 \operatorname{Haf}X=\sum_{M\in\mathcal M_{2n}}\prod_{\{i,j\}\in M}X_{ij},
 \qquad h_n=(2n-1)!!=\frac{(2n)!}{2^n n!}.
 \label{eq:ghe-hafnian}
\end{equation}
Given $0<\varepsilon,\delta<1$, an estimator must output $\widehat p\in\mathbb R$ with
\begin{equation}
 \Pr\!\left[\left|\widehat p-|\operatorname{Haf}X|^2\right|
                   \leq\varepsilon h_n\right]\geq1-\delta.
 \label{eq:ghe-estimation}
\end{equation}
Probability in Eq.~\eqref{eq:ghe-estimation} includes both $X$ and the estimator's randomness.
The binary input is a truncated, rounded copy $\widetilde X$.
Use $\operatorname{poly}(n,\log(1/\varepsilon),\log(1/\delta))$ bits per entry.
Choose the precision so that the squared hafnians of $X$ and $\widetilde X$ differ by at most $\varepsilon h_n/2$, except with probability $\delta/2$.
Reduction cost is polynomial in the binary input length, $1/\varepsilon$, and $1/\delta$.
Here $\#\mathrm P$ is the class of functions counting accepting paths of nondeterministic polynomial-time machines.
The target normalization $h_n$ is defined in Eq.~\eqref{eq:ghe-hafnian}.

\subsection*{Status}

Unsolved

\subsection*{Source}

This explicit ensemble and additive normalization refine the hafnian-of-Gaussians discussion following Eq.~(11) of Hamilton et al. \sourcecite{ref:ghe-hamilton}{HKS+17}.
The formulation distinguishes additive squared-hafnian estimation from multiplicative hafnian estimation; the paper does not state this precise normalized problem as a numbered conjecture.

\subsection*{Progress}

\begin{itemize}
  \item Hamilton et al., Eq.~(7), use the identity
\begin{equation}
 \operatorname{Haf}\begin{pmatrix}0&G\\G^T&0\end{pmatrix}
 =\operatorname{Per}G.
 \label{eq:ghe-block}
\end{equation}
Here $G$ is any square matrix and $\operatorname{Per}G$ is its permanent.
Equation~\eqref{eq:ghe-block} transfers worst-case permanent evaluation hardness to hafnians.
The block matrices in this reduction are not typical independent symmetric Gaussian samples \sourcecite{ref:ghe-hamilton}{HKS+17}.

  \item Independence and the matching expansion give $\mathbb E|\operatorname{Haf}X|^2=h_n$.
Distinct matchings have vanishing cross terms.
This elementary normalization calculation does not provide an average-case reduction at error $\varepsilon h_n$.
\end{itemize}

\subsection*{References}

\begin{enumerate}[label={},leftmargin=4.8em,labelsep=0.5em]
  \footnotesize
  \item[\textup{[HKS+17]}]\label{ref:ghe-hamilton}
  C. S. Hamilton, R. Kruse, L. Sansoni, S. Barkhofen, C. Silberhorn, and I. Jex, "Gaussian Boson Sampling," \emph{Physical Review Letters} \textbf{119}, 170501 (2017). \href{https://doi.org/10.1103/PhysRevLett.119.170501}{doi:10.1103/PhysRevLett.119.170501}; \href{https://arxiv.org/abs/1612.01199}{arXiv:1612.01199}.
\end{enumerate}

\subsection*{Comment}

The unresolved claim is average-case additive approximation hardness for independent complex entries above the diagonal.
Worst-case hardness, exact average-case evaluation, and Gaussian-product ensembles $YY^T$ do not by themselves imply it.
The hafnian lower-tail question is a separate auxiliary problem, not an equivalent hardness claim.

\subsection*{Field}

Quantum algorithm

\subsection*{Topic}

Boson sampling; Computational complexity and computability

\subsection*{ID}

\texttt{op\_903c729ed6eb9ce6}
